\section{Hur jag ska förbättra mina kunskaper och färdigheter}
\label{sec:forbattra_kunskaper}

LIA 2 har hjälpt mig att se både vad jag kan och vad jag fortfarande behöver utveckla. Jag tycker
att det är viktigt att vara tydlig med båda delarna, eftersom det gör det lättare att sätta rätt mål
för nästa steg.

\subsection*{Observability -- det jag saknar mest}

Under LIA 2 arbetade jag mycket med att bygga och deploya system. I efterhand ser jag att jag
vill bli bättre på att övervaka system efter att de är driftsatta. Vi hade loggar, men inte ett
fullt övervakningssystem med tydliga dashboards och larm.

Det här är viktigt eftersom drift inte bara handlar om att starta ett system, utan också om att
tidigt se när något börjar gå fel. Därför vill jag lära mig mer om verktyg som Prometheus,
Grafana och centraliserad loggning.

\subsection*{Deployment patterns -- Blue/Green och Canary}

Jag löste ett konkret problem med driftstopp under deployment, men jag vill nu lära mig mer om
andra sätt att rulla ut ny kod, till exempel Blue/Green deployment och Canary releases. Jag
tror att den kunskapen blir viktig när jag senare arbetar i större miljöer med högre krav på
stabilitet.

\subsection*{Terraform i större skala}

Under LIA 2 använde jag Terraform CDKTF i ett konkret projekt, men jag vill bli bättre på hur
man strukturerar Terraform i större skala. Jag vill förstå mer om moduler, state-hantering och
hur man arbetar med flera miljöer på ett säkert sätt.

Det är ett område där jag känner att jag har en bra grund, men där jag fortfarande behöver mer
erfarenhet.

\subsection*{Certifieringarna som struktur för lärandet}

Jag vill använda certifieringar som ett sätt att fortsätta utvecklas på ett strukturerat sätt.
Planen är att börja med CKA efter examen och därefter fortsätta med Google Professional Cloud
DevOps Engineer. På det sättet kan jag kombinera praktisk erfarenhet med mer systematisk
fördjupning.